% Generated draft: requires mathematical and source review.
\documentclass[11pt]{article}
\usepackage[margin=25mm]{geometry}
\usepackage{fontspec}
\setmainfont{Latin Modern Roman}
\usepackage{amsmath,amssymb,mathtools,unicode-math}
\setmathfont{Latin Modern Math}
\usepackage{enumitem,needspace,xurl,hyperref}
\setlength{\parindent}{0pt}
\setlength{\parskip}{5pt}
\setlength{\emergencystretch}{3em}
\newcommand{\N}{\mathbb{N}}
\newcommand{\Z}{\mathbb{Z}}
\newcommand{\Q}{\mathbb{Q}}
\newcommand{\R}{\mathbb{R}}
\newcommand{\C}{\mathbb{C}}
\newcommand{\F}{\mathbb{F}}
\newcommand{\A}{\mathbb{A}}
\newcommand{\PP}{\mathbb{P}}
\newcommand{\E}{\mathbb{E}}
\newcommand{\eps}{\varepsilon}
\newcommand{\dd}{\,\mathrm{d}}
\begin{document}
\section*{UnsolvedMath ID: 10400133}
\textbf{Source problem code:} catalogue entry
\subsection{Definitions and mathematical statement}
A flat principal \(G\)-bundle over a manifold \(M\) is a principal bundle equipped with a flat connection; for connected \(M\), its holonomy is represented by a homomorphism \(\pi_1(M)\to G\), up to the appropriate conjugacy equivalence.

S. Baseilhac and R. Benedetti ask to generalize the construction of the QHI, the source's quantum hyperbolic invariants, from flat principal \(B\)-bundles to flat principal \(G\)-bundles for Lie groups \(G\ne B\). The object to be generalized includes dependence on the flat bundle, not merely on the underlying manifold or the abstract group.

The symbol \(B\) is not defined in the excerpt. In the usual quantum-hyperbolic setting it denotes the upper triangular subgroup
\[
 \left\{\begin{pmatrix}a&b\\0&a^{-1}\end{pmatrix}:a\in\C^\times,\ b\in\C\right\}\subset\operatorname{SL}_2(\C),
\]
but that identification is explicitly provisional. The original QHI formula, admissible manifolds, quantum parameters, and any link, triangulation or decoration data are also absent. A valid extension must specify its full input and prove it is well defined under the corresponding equivalences, with a clear relation to the original construction. No additional auxiliary data are silently discarded.

The source does not identify which other Lie groups, whether one uniform construction for every such group is requested, or which structural properties of QHI must persist. Thus the generalization remains a construction programme rather than a specified formula for a universal invariant.
\subsection{Short English statement}
Extend the source's quantum hyperbolic invariant construction to flat bundles with structure groups beyond \(B\), retaining bundle dependence and necessary auxiliary data. The original construction and intended other groups are not specified in the extract.
\subsection{Sources}
\begin{itemize}
\item Local UnsolvedMath registry, permanent ID 10400133, source code catalogue entry; statement attribution S. Baseilhac and R. Benedetti; proposer field and proposal year absent; updated 2026-08-06.
\item Ohtsuki, \emph{Problems on invariants of knots and 3-manifolds} (2002), Problem 7.18, PDF page 113: \url{https://msp.org/gtm/2002/04/gtm-2002-04-024s.pdf}. The excerpt explicitly credits Section 7.4 to Baseilhac and Benedetti. Registry access date 2026-07-29; PDF font-structure extraction; provenance distinguishes 231 numbered bold headings. Present status and redistribution terms require review.
\item The supplied summary repeats the QHI generalization question without an exact additional construction reference or verified resolution.
\end{itemize}
\end{document}
